\chapter{自动化混音}
\label{chap:automation}

\classquote{穷则变，变则通，通则久。}{《周易·系辞下》}{简注：走到尽头就要变，变了才通，通了才久；自动化是在时间轴上随段落改变推子与参数，而不是一个静态开关。}

静态混音回答“整首歌用同一个推子位置”，自动化回答“第几小节推到哪里”。夜半的自动化不是插件里的一个“自动化开关”，而是挂在每条音轨上的一组曲线：每个点记录一个 tick 和一个值，播放或导出时由模型按曲线求值，再用这个值取代（不是叠加）该目标的静态值。本章只写活尖端真实存在的目标、泳道、读/写标志与编辑入口；没有做出来的部分在末尾逐条列出。

\section{自动化在夜半的层级}

一条自动化泳道由三样东西决定：目标、点、读/写状态。目标回答“控制谁”，点回答“什么时候是多少”，读/写状态回答“播放时用不用、录制时写不写”。

\begin{table}[htbp]
\centering
\small
\begin{tabular}{@{}p{0.18\linewidth}p{0.26\linewidth}p{0.18\linewidth}p{0.26\linewidth}@{}}
\hline
目标变体 & 寻址 & 单位与标称显示域 & 没有自动化时用什么 \\
\hline
\texttt{TrackVolume} & 音轨 ID & dB，$-60$–$+12$ & 音轨静态 \texttt{volume\_db} \\
\texttt{TrackPan} & 音轨 ID & 双极，$-1$–$+1$ & 音轨静态 \texttt{pan} \\
\texttt{SendGain} & 音轨 ID ＋ 路由边 ID & dB，$-60$–$+12$ & 该路由边的 \texttt{gain\_db}；缺省为单位增益 0 dB \\
\texttt{DeviceParam} & 音轨 ID ＋ 设备槽位 ＋ 参数下标 & 参数原生单位；模型不声明值域 & 该设备参数的当前值 \\
\texttt{Macro} & 音轨 ID ＋ 宏下标 & 归一化，$0$–$1$ & 宏当前位置 \\
\texttt{MidiCc} & 音轨 ID ＋ 控制器号 & 原始 7 位值，$0$–$127$ & 该控制器缺省 0 \\
\hline
\end{tabular}
\caption{夜半的六种自动化目标；本章重点写前四种，宏与 MIDI CC 只列入口}
\label{tab:automation-targets}
\end{table}

“标称显示域”只是画纵轴时的缺省范围，不是写入钳位：模型只要求点的值有限，曲线可以落到显示域之外，界面只把像素钳进泳道带。不要把 $-60$–$+12$ 当成“推子超出这个范围就会被拒绝”。

\begin{center}
\begin{tikzpicture}[
  box/.style={draw, rounded corners=2pt, minimum width=30mm, minimum height=9mm, align=center, font=\scriptsize},
  arrow/.style={->, thick}
]
  \node[box] (sec) at (0,0) {段落/小节\\（编辑者的时间分区）};
  \node[box] (pts) at (4.0,0) {泳道上的点\\tick + 值 + 曲线};
  \node[box] (eval) at (8.4,0) {模型求值\\唯一入口};
  \node[box] (live) at (6.2,-2.0) {实时监听\\音量/声相；部分发送};
  \node[box] (off) at (10.3,-2.0) {离线导出\\全部可求值目标};
  \draw[arrow] (sec) -- (pts);
  \draw[arrow] (pts) -- (eval);
  \draw[arrow] (eval) -- (live);
  \draw[arrow] (eval) -- (off);
\end{tikzpicture}
\end{center}

段落不是自动化的容器。夜半有曲式段落实体，但自动化泳道只认音轨上的 tick；移动、复制或改段落不会自动搬走曲线。把“主歌轻、副歌推起来”写进曲线，仍然要按小节位置自己落点。

\section{泳道怎么显示}

编排视图里，每条轨可以有多条自动化泳道；同一条轨的多条泳道按目标类型的固定顺序等分车道行的高度。一个目标一条带，带与带之间不会重叠。默认是展开的；\texttt{Cmd/\allowbreak Ctrl+Shift+A} 或“视图 → 显示 / 隐藏自动化泳道”切换显示。菜单文案写“(A)”，实际绑定是带 \texttt{Shift} 的组合，单按 \texttt{A} 不会切换。

\begin{table}[htbp]
\centering
\small
\begin{tabular}{@{}p{0.24\linewidth}p{0.60\linewidth}@{}}
\hline
界面元素 & 事实 \\
\hline
泳道标签 & 音量显示“音量”，声相显示“声相”，发送显示“发送 <路由边 ID>”，设备参数显示参数名，宏显示宏名，MIDI CC 显示“CC <号>” \\
读/写角标 & 读关闭时显示“读关”；写模式非 \texttt{Off} 时显示“$\bullet$ 写入/触碰/锁存”；两者可同时出现 \\
曲线颜色 & 读打开时用亮色；读关闭时用暗色，但点仍在 \\
主总线 & 主总线轨的泳道不在编排视图绘制；主总线的混音在调音台通道条上做 \\
空泳道 & 没有点的泳道不持久化；界面只画工程里已经存在的泳道 \\
\hline
\end{tabular}
\caption{自动化泳道的显示口径}
\label{tab:automation-display}
\end{table}

\screenshot{automation-lane-overview}{编排视图展开多条自动化泳道，可见音量、声相与设备参数带及读/写角标}

点击泳道带会在该位置增/改一个点：横坐标换算成 tick，纵坐标按该泳道的有效显示域换算成值，新增点的曲线缺省为线性。已经有点落在点击位置 8 像素以内时，改的是既有点的值，保留它原来的曲线和身份；不会堆出第二个同 tick 的点。这个点击入口只能改“已存在泳道”的点，不能新建泳道，也不能拖动点、删点或改曲线形状。

\begin{table}[htbp]
\centering
\small
\begin{tabular}{@{}p{0.20\linewidth}p{0.24\linewidth}p{0.46\linewidth}@{}}
\hline
操作 & 入口 & 边界 \\
\hline
显示/隐藏 & \texttt{Cmd/\allowbreak Ctrl+Shift+A}；视图菜单 & 默认展开；没有“只显示选中轨”的过滤 \\
增/改点 & 在泳道带内点击 & 只能改已存在泳道；8 像素内算同一点；新点为线性 \\
删点 & 无界面入口 & 走 MCP 的 \texttt{removeAutomationPoint} \\
改曲线 & 无界面入口 & 走 MCP 写点的 \texttt{curve}，或工程 JSON \\
新建泳道 & 无界面入口 & 先写第一个点（MCP \texttt{yeban\_edit\_automation}）或改工程 JSON \\
读/写开关 & 无界面入口 & 角标只显示；改属性走 MCP 的 \texttt{set\allowbreak Automation\allowbreak Lane} \\
\hline
\end{tabular}
\caption{自动化点的界面编辑入口}
\label{tab:automation-edit-entry}
\end{table}

\screenshot{automation-point-edit}{在泳道带内点击后新增的自动化点，曲线与读/写角标同时可见}

\section{点、曲线与求值}

夜半用 960 PPQ 的整数 tick，不用采样数或浮点秒。常用换算：

\begin{table}[htbp]
\centering
\small
\begin{tabular}{@{}lll@{}}
\hline
时值 & tick & 4/4 下的位置 \\
\hline
四分音符（1 拍） & 960 & 1 小节 = 4 拍 = 3{,}840 tick \\
十六分音符 & 240 & 1 小节 = 16 个十六分 \\
1 小节 & 3{,}840 & 8 小节 = 30{,}720 tick \\
\hline
\end{tabular}
\caption{自动化点的 tick 换算}
\label{tab:automation-ticks}
\end{table}

相邻两个点之间按左端点的曲线形状插值；曲线只改变“到下一个点”的形状，不改变点本身的值。四种曲线都在端点精确落在值上，没有过冲：

\begin{table}[htbp]
\centering
\small
\begin{tabular}{@{}lll@{}}
\hline
序列化名 & 归一化权重 $u(t)$ & 听感/用途 \\
\hline
\texttt{Linear} & $t$ & 等速推子；最常用 \\
\texttt{Exponential} & $t^2$ & 先慢后快；适合“最后再推起来” \\
\texttt{Logarithmic} & $t(2-t)$ & 先快后慢；适合“先推到位再稳住” \\
\texttt{SCurve} & $t^2(3-2t)$ & 两端缓、中间快；适合段落交接 \\
\hline
\end{tabular}
\caption{自动化曲线；$t$ 是相邻两点之间的归一化位置}
\label{tab:automation-curves}
\end{table}

边界规则：

\begin{itemize}
  \item 第一个点之前保持第一个点的值；最后一个点之后保持最后一个点的值。没有外推。
  \item 空泳道没有自动化值，播放/导出退回静态值。
  \item 同一个 tick 上有多个点时，按点 ID 的字典序取最大者；界面点击会更新既有点，正常操作不会制造同 tick 多点。
  \item 值超出显示域时，模型保留原值，界面只把像素钳进带内；不要用“曲线贴顶”反推“值被钳到上界”。
\end{itemize}

\section{读开关与写模式}

每条泳道有两个独立状态：\texttt{readEnabled} 与 \texttt{writeMode}。读开关缺省为真；关掉后，求值入口直接返回“无自动化值”，播放与导出都退回静态值，曲线还在，但不再生效。写模式缺省为 \texttt{Off}；它不是“录制按钮”，只是数据里的一个门控。

\begin{table}[htbp]
\centering
\small
\begin{tabular}{@{}p{0.14\linewidth}p{0.30\linewidth}p{0.44\linewidth}@{}}
\hline
写模式 & 设计语义 & 当前实现 \\
\hline
\texttt{Off} & 不录制 & 门控关闭，任何手势都不写点 \\
\texttt{Write} & 走带时只要参数被改动就写入 & 只在音量/声相/EQ 低架手势松手时写一个点；没有“走带即写” \\
\texttt{Touch} & 仅手动触碰时写入，松手即停 & 与 \texttt{Write} 走同一条松手写入路径；没有额外分流 \\
\texttt{Latch} & 一直写到显式停止 & 与 \texttt{Write} 走同一条松手写入路径；没有“锁存到停止” \\
\hline
\end{tabular}
\caption{四种写模式的设计语义与当前实现；三种非 \texttt{Off} 模式尚未分流}
\label{tab:automation-write-modes}
\end{table}

界面目前只把读/写状态画成角标，没有开关、录制臂或模式选择器。要改 \texttt{readEnabled}、\texttt{writeMode} 或显式 \texttt{domain}，走 MCP 的 \texttt{yeban\_edit\_notes}，在 \texttt{ops} 里给 \texttt{kind: "set\allowbreak Automation\allowbreak Lane"}；把整条泳道取走用同一条里的 \texttt{remove: true}。新建泳道不能先建一条空泳道：模型拒绝持久化“无点、读开、Off、无覆盖”的隐式形状；先用 \texttt{yeban\_edit\_automation} 写第一个点，泳道随点出现，再改属性。

\subsection{录制实际只写一个点}

自动化的“录制”只在播放中、且手势松手时发生。当前写点的手势只有三类：

\begin{enumerate}
  \item 混音台轨道推子（\texttt{TrackVolume}）；
  \item 混音台轨道声相（\texttt{TrackPan}）；
  \item 设备机架的 EQ 低架增益（\texttt{DeviceParam}，只针对 EQ 的低架增益参数）。
\end{enumerate}

松手时取当前播放头的 tick，写一个点；拖动过程中的中间值不会形成曲线。按下就松手、拖回原值、播放停止、泳道不存在、写模式为 \texttt{Off}，五种情况都不写。发送增益推子、宏、MIDI CC 没有自动录制路径；发送增益的静态推子只改路由边增益，不写 \texttt{SendGain} 点。

因此“写模式”在当前版本的实际作用是“这条泳道允不允许松手写一个点”，不是四种不同的录制行为。需要连续曲线时，只能点击落点或走 MCP 批量写点。

\section{四类目标怎么写}

\subsection{音量}

\texttt{TrackVolume} 的值是绝对 dB，不是相对推子的偏移。存在自动化时，静态 \texttt{volume\_db} 被曲线值取代；没有自动化时才用静态值。想让整首在 $-6$ dB 附近，曲线也写 $-6$，不要写 $0$ 再指望它叠加。

段落用法：主歌压低、副歌还原、尾奏收。示例（8 小节，4/4，30{,}720 tick）：

\begin{table}[htbp]
\centering
\small
\begin{tabular}{@{}llll@{}}
\hline
tick & 小节 & 值 & 意图 \\
\hline
0 & 1 & $-18$ dB & 前奏从很低开始 \\
7{,}680 & 3 & $-6$ dB & 主歌进来 \\
15{,}360 & 5 & $0$ dB & 副歌满推子 \\
23{,}040 & 7 & $0$ dB & 保持 \\
30{,}720 & 9（尾） & $-12$ dB & 尾奏收掉 \\
\hline
\end{tabular}
\caption{音量自动化示例；值为绝对 dB}
\label{tab:automation-volume-recipe}
\end{table}

\subsection{声相}

\texttt{TrackPan} 的值是 $-1$–$+1$ 的绝对位置。它走第 7 章的声相衰减律与“中心恒等”规则：读数为 $0$ 时左右都不乘增益；一旦离开中心，立刻落到所选律的曲线。默认 $-3$ dB 律在中心与极小非零之间有 $3.01$ dB 台阶；自动化从 $0$ 慢慢移到 $0.6$ 时，离开中心的那一刻会有这个台阶，不能用它做“无缝从中间展开”。要避免台阶，段落内要么保持中心，要么直接从非零位置开始，别让曲线穿过 $0$。

示例：主歌保持中心，副歌把双声道素材展开到 $\pm0.6$，尾奏回中：

\begin{table}[htbp]
\centering
\small
\begin{tabular}{@{}llll@{}}
\hline
tick & 小节 & 值 & 意图 \\
\hline
0 & 1 & $0$ & 前奏居中 \\
7{,}680 & 3 & $0$ & 主歌仍居中 \\
15{,}360 & 5 & $+0.6$ & 副歌向右展开 \\
23{,}040 & 7 & $+0.6$ & 保持 \\
30{,}720 & 9（尾） & $0$ & 回中 \\
\hline
\end{tabular}
\caption{声相自动化示例；副歌展开，尾奏回中}
\label{tab:automation-pan-recipe}
\end{table}

如果素材本身是单声道，声相自动化只改变左右分配，不会产生新的立体声信息；宽度仍走通道条的 M/S 级（第 7 章）。

\subsection{发送}

\texttt{SendGain} 寻址的是路由边，不是“第几个发送”。界面发送行不显示边 ID；泳道标签显示“发送 <边 ID>”，要从工程查询或工程 JSON 里取 \texttt{edgeId}。值是该边的绝对增益 dB，不是相对静态发送量的偏移；路由边 \texttt{gain\_db} 为 \texttt{None} 时静态值是单位增益 0 dB。静态发送量是 $-6$ dB 时，曲线第一个点也写 $-6$，否则一开自动化就会把发送量跳到 0 dB。

示例：主歌保持 $-6$ dB，副歌第一拍做一个 $+3$ dB 的短推，再回 $-6$：

\begin{table}[htbp]
\centering
\small
\begin{tabular}{@{}llll@{}}
\hline
tick & 小节 & 值 & 意图 \\
\hline
0 & 1 & $-6$ dB & 与静态发送量对齐 \\
15{,}360 & 5 & $-6$ dB & 副歌起点仍保持 \\
17{,}280 & 5.5 & $-3$ dB & 短推 \\
19{,}200 & 6 & $-6$ dB & 回到原发送量 \\
30{,}720 & 9（尾） & $-6$ dB & 保持 \\
\hline
\end{tabular}
\caption{发送增益自动化示例；值是绝对 dB}
\label{tab:automation-send-recipe}
\end{table}

发送增益推子本身不写 \texttt{SendGain} 点；上面这条曲线要靠点击泳道或 MCP 写。发送泳道的实时落点取决于边是否真的搬运音频：有落点的边在监听里生效；查不到、侧链、源是母线或源不是可渲染音轨的边，只在离线导出里生效。

\subsection{效果}

\texttt{DeviceParam} 寻址“第几个设备槽的第几个参数”。值是参数的原生单位，模型不声明通用值域；纵轴要么用泳道显式 \texttt{domain}，要么按曲线最值自适应。界面只有 EQ 低架增益有可拖动的参数自动化入口；其他设备参数只能走 MCP 写点。

示例：EQ 低架增益（参数域 $-18$–$+18$ dB，缺省 0 dB）在主歌略减、副歌略增：

\begin{table}[htbp]
\centering
\small
\begin{tabular}{@{}llll@{}}
\hline
tick & 小节 & 值 & 意图 \\
\hline
0 & 1 & $-2$ dB & 主歌略薄 \\
15{,}360 & 5 & $-2$ dB & 副歌前仍保持 \\
17{,}280 & 5.5 & $+3$ dB & 副歌抬低架 \\
23{,}040 & 7 & $+3$ dB & 保持 \\
30{,}720 & 9（尾） & $-2$ dB & 回主歌值 \\
\hline
\end{tabular}
\caption{EQ 低架增益自动化示例；参数原生单位 dB}
\label{tab:automation-device-recipe}
\end{table}

\screenshot{automation-device-param}{设备机架里的 EQ 低架增益与对应的设备参数泳道，二者共用同一个槽位下标}

设备参数泳道存的是槽位下标，不是设备身份。若在目标设备前面插入或删除一台设备，同一个槽位会指向另一台设备的同下标参数；夜半没有按 \texttt{device\_id} 的稳定寻址。改设备链后必须重新核对泳道落点。

\section{按段落设计曲线}

把“段落”翻译成 tick 区间，再决定每个区间用哪条目标，不要把所有变化都堆在音量上：

\begin{table}[htbp]
\centering
\small
\begin{tabular}{@{}p{0.14\linewidth}p{0.18\linewidth}p{0.24\linewidth}p{0.30\linewidth}@{}}
\hline
段落 & tick 区间 & 主要目标 & 写法 \\
\hline
前奏 & 0–7{,}680 & 音量 & 从 $-18$ 到 $-6$，线性或 S 曲线 \\
主歌 & 7{,}680–15{,}360 & 音量、声相 & 音量 $-6$；声相保持中心 \\
副歌 & 15{,}360–23{,}040 & 音量、发送、效果 & 音量 0；发送短推；EQ 低架略增 \\
尾奏 & 23{,}040–30{,}720 & 音量、声相 & 音量回 $-12$；声相回中 \\
\hline
\end{tabular}
\caption{8 小节段落的自动化分配}
\label{tab:automation-section-plan}
\end{table}

操作顺序：

\begin{enumerate}
  \item 先定静态值：推子、声相、发送量、设备参数都先摆到“没有自动化时”的基准。
  \item 要新建泳道时，先用 MCP 写第一个点，让泳道出现；再用 \texttt{set\allowbreak Automation\allowbreak Lane} 设定读/写与显示域。
  \item 在编排视图按 \texttt{Cmd/\allowbreak Ctrl+Shift+A} 确认泳道可见；点击落点，或走 MCP 批量写点。
  \item 需要录制时，先确认泳道 \texttt{writeMode} 不是 \texttt{Off}，再播放，拖推子/声相/EQ 低架，松手写一个点。
  \item 读回：用 \texttt{yeban\_edit\_automation} 的 \texttt{ticks} 查“自动化值/静态值/实际值”，确认 \texttt{source} 是 \texttt{automation} 还是 \texttt{static}。
  \item 导出前确认读开关没有关；关掉的泳道在导出里同样不生效。
\end{enumerate}

实时监听与离线导出的落点不同，必须分开看：

\begin{table}[htbp]
\centering
\small
\begin{tabular}{@{}p{0.22\linewidth}p{0.26\linewidth}p{0.42\linewidth}@{}}
\hline
目标 & 实时监听 & 离线导出 \\
\hline
\texttt{TrackVolume} & 生效 & 生效 \\
\texttt{TrackPan} & 生效 & 生效 \\
\texttt{SendGain} & 边能搬运音频时生效 & 生效；边无实时落点时只在导出里生效 \\
\texttt{DeviceParam} & 只有通道条来源的 \texttt{output\_gain\_db} 生效 & 插入链标量参数生效；未知名字不生效 \\
\texttt{Macro} & 不生效 & 不生效 \\
\texttt{MidiCc} & 不生效 & 生效 \\
\hline
\end{tabular}
\caption{自动化在监听与导出里的落点}
\label{tab:automation-apply}
\end{table}

这意味着“在监听里拧 EQ 低架并录了一个点”并不等于“监听里听到了 EQ 变化”：EQ 低架属于 \texttt{DeviceParam}，实时侧没有该参数的槽位，点会写进工程，但监听不执行，导出才执行。混音时不要用监听判断这条曲线的效果，至少导出一次核对。

\section{尚未提供与边界}

\begin{table}[htbp]
\centering
\small
\begin{tabular}{@{}p{0.36\linewidth}p{0.54\linewidth}@{}}
\hline
需求 & 夜半现状 \\
\hline
在界面上新建/删除泳道 & 只能先写点或改工程；没有“添加自动化目标”的菜单 \\
拖点、框选多点、批量删除 & 没有；点击只能增/改单点 \\
改点曲线形状 & 没有；点击保留原曲线，新点固定线性 \\
读/写开关与模式选择器 & 没有；角标只读，改属性走 MCP \\
\texttt{Write}/\texttt{Touch}/\texttt{Latch} 分流 & 没有；三种非 \texttt{Off} 模式都只走松手写一点 \\
连续录制一条推子曲线 & 没有；一次松手只写一个点 \\
发送增益推子录制 & 没有；只改路由边静态增益 \\
宏/MIDI CC 录制 & 没有 \\
段落级自动化容器 & 没有；曲线按整条轨的 tick 轴，段落移动不带曲线 \\
主总线泳道可视化 & 编排视图不画主总线泳道 \\
设备参数的稳定身份寻址 & 没有；泳道存槽位下标，设备链变动可能漂移 \\
自动化点稀疏化/简化 & 没有；点写多少就存多少 \\
\hline
\end{tabular}
\caption{自动化混音的当前边界}
\label{tab:automation-gaps}
\end{table}

\section{练习}

\begin{enumerate}
  \item 给一条轨画 4 小节的音量自动化：第一小节静态、第二小节渐强、第三小节渐弱、第四小节回静态；回放确认曲线与听感一致。
  \item 比较静态推子与自动化在同一段的差别；记录自动化是否真的在改变电平，而不是只改了显示。
  \item 查一次读/写开关：录制一个点后停止，确认没有把后续手动推子写成新点。
\end{enumerate}

\section{本章与帮助主题的映射}

\begin{table}[htbp]
\centering
\small
\begin{tabular}{@{}p{0.34\linewidth}p{0.52\linewidth}@{}}
\hline
本章小节 & 应用内帮助锚点 \\
\hline
自动化在夜半的层级 & \texttt{automation-layers} \\
泳道怎么显示 & \texttt{automation-lanes} \\
点、曲线与求值 & \texttt{automation-points-curves} \\
读开关与写模式 & \texttt{automation-read-write} \\
录制实际只写一个点 & \texttt{automation-record-one-point} \\
音量自动化 & \texttt{automation-volume} \\
声相自动化 & \texttt{automation-pan} \\
发送自动化 & \texttt{automation-send} \\
效果自动化 & \texttt{automation-device-param} \\
段落设计与落点 & \texttt{automation-section-plan} \\
当前边界 & \texttt{automation-gaps} \\
\hline
\end{tabular}
\caption{第八章帮助主题映射}
\label{tab:automation-help-map}
\end{table}

\section*{延伸阅读}
\begin{itemize}[leftmargin=1.5em]
\item \url{https://www.ableton.com/en/blog/}——Ableton 博客：看自动化泳道与段落变化的制作案例，补足“何时动、动多少”的判断。
\item \url{https://github.com/gradetwo/yeban}——夜半仓库：按源码核对自动化目标、泳道与读/写标志的真实实现边界。
\end{itemize}
